The experimental results validate the application of Exponential Fourier Series (EFS) in decomposing periodic signals and characterizing continuous-time systems. In the first phase, the reconstruction of the sawtooth signal $x(t) = -2t$ using calculated Fourier coefficients $C_n$ demonstrated that as the number of terms $n$ increases, the approximation converges toward the original function. However, the presence of oscillations near the points of discontinuity confirms the existence of the Gibbs phenomenon. As noted in [1], this phenomenon is an inherent mathematical limitation where a finite sum of continuous sine and cosine waves cannot perfectly represent a jump discontinuity, resulting in a persistent overshoot of approximately 9\%. This necessitates the consideration of the Dirichlet conditions when determining the representability of a signal in the frequency domain [2].

The second phase of the laboratory focused on the system impulse response $h(t) = \frac{a}{2}e^{-a|t|}$. The data revealed a linear proportional relationship between the exponential scaling factor $a$ and the 3dB cut-off frequency $\omega_c$. As the parameter $a$ increased from 10 to 100, the system's magnitude response $|H(\omega)|$ exhibited a widening passband. This behavior identifies the system as a low-pass filter where $a$ serves as the primary control for bandwidth expansion. Such parametric control is vital in practical engineering applications, as it allows for the precise adjustment of a filter’s spectral envelope to either isolate signals of interest or suppress high-frequency noise [3].

In conclusion, the experiment successfully demonstrated the transition from time-domain signal modeling to frequency-domain system analysis. The use of MATLAB's Symbolic Math Toolbox allowed for the accurate calculation of $C_n$ and the visualization of parametric shifts in filter performance. The findings underscore the trade-offs between series truncation and signal fidelity, as well as the direct impact of system parameters on absolute bandwidth. These principles are fundamental to the design of communication systems and signal processing chains where spectral efficiency and reconstruction accuracy are paramount [4].
